\section{RELATÓRIO DE TRABALHO --- SITE INSTITUCIONAL
CDA}\label{relatuxf3rio-de-trabalho-site-institucional-cda}

\textbf{Data:} 07 de Outubro de 2026\\
\textbf{Elaborado para:} Direcção da CDA\\
\textbf{Projecto:} Site Institucional da Câmara dos Despachantes
Aduaneiros de Moçambique (CDA)

\subsection{1. INTRODUÇÃO}\label{introduuxe7uxe3o}

Este relatório apresenta, de forma clara e objectiva, o estado actual do
projecto de redesenho e desenvolvimento do Site Institucional da CDA.
Descreve o que já foi realizado desde o início, o que está em andamento
e o que ainda falta executar, com indicação de prioridades.

\subsection{2. OBJECTIVO DO TRABALHO}\label{objectivo-do-trabalho}

Desenvolver um site institucional moderno, responsivo, acessível e
profissional para a CDA, alinhado com a sua identidade corporativa, com
estrutura clara de informação, dados organizados e bom desempenho.

\subsection{3. RESUMO EXECUTIVO}\label{resumo-executivo}

\textbf{Estado Geral:} Em desenvolvimento --- Fundação sólida concluída
(70--75\% da estrutura base).

\begin{itemize}
\tightlist
\item
  \textbf{Concluído (Base Estrutural):} Design System com cores CDA,
  tipografia profissional, dados institucionais organizados e validados,
  componentes interativos funcionais.
\item
  \textbf{Em Andamento:} Aplicação do novo design às páginas do site.
\item
  \textbf{Por Concluir:} Migração completa das restantes páginas para o
  novo padrão visual.
\end{itemize}

\subsection{4. O QUE FOI FEITO
(Realizado)}\label{o-que-foi-feito-realizado}

\subsubsection{4.1 Identidade Visual e Design
System}\label{identidade-visual-e-design-system}

\begin{itemize}
\tightlist
\item
  Criado Design System próprio para o projecto (tokens, sistema e
  layout).
\item
  Definida paleta oficial com cores da CDA (\#B71C1C --- Vermelho
  Institucional, \#0F3D56 --- Azul Naval).
\item
  Implementada tipografia \textbf{Manrope} (variável 400--800),
  auto-hospedada (sem dependência externa).
\item
  Garantida conformidade com requisitos de contraste (WCAG AA) ---
  verificado por script automatizado.
\item
  Criados estilos para componentes reutilizáveis (botões, cards, avisos,
  separadores/tabs).
\end{itemize}

\subsubsection{4.2 Estrutura, Navegação e
Interface}\label{estrutura-navegauxe7uxe3o-e-interface}

\begin{itemize}
\tightlist
\item
  Desenvolvido novo cabeçalho institucional com mega-menu organizado por
  secções.
\item
  Desenvolvido rodapé institucional completo com contactos e links
  organizados.
\item
  Implementada caixa de pesquisa global com acessibilidade (teclado,
  leitor de ecrã, ESC).
\item
  Adicionadas funcionalidades de animação/reveal com bom desempenho.
\item
  Componentes com foco em acessibilidade (ARIA, navegação por teclado).
\end{itemize}

\subsubsection{4.3 Dados Institucionais (Base de Dados do
Site)}\label{dados-institucionais-base-de-dados-do-site}

\begin{itemize}
\tightlist
\item
  \textbf{Membros/Despachantes:} 218 registos validados, sem dados
  sensíveis (PII).
\item
  \textbf{Eventos:} 38 registos organizados e categorizados.
\item
  \textbf{Cronologia/Historial:} 19 registos com base em documentação
  oficial da CDA.
\item
  \textbf{Documentação/PDFs:} 266 documentos mapeados e organizados.
\item
  Todos os conjuntos de dados validados automaticamente (0 erros).
\end{itemize}

\subsubsection{4.4 Funcionalidades
Desenvolvidas}\label{funcionalidades-desenvolvidas}

\begin{itemize}
\tightlist
\item
  \textbf{Directório de Membros}: Listagem, pesquisa, filtros e
  paginação --- testado e funcional.
\item
  \textbf{Cronologia Institucional}: Apresentação cronológica dos factos
  mais relevantes --- funcional.
\item
  \textbf{Galeria de Imagens}: Mosaico responsivo com lightbox acessível
  --- funcional.
\item
  \textbf{Separadores (Tabs)}: Componentes reutilizáveis para
  organização de conteúdos --- funcionais.
\item
  \textbf{Lógica de dados unificada}: Garantida uma única fonte de dados
  para evitar inconsistências.
\end{itemize}

\subsubsection{4.5 Qualidade, Testes e
Documentação}\label{qualidade-testes-e-documentauxe7uxe3o}

\begin{itemize}
\tightlist
\item
  Testes de links: \textbf{0 falhas} detectadas (208 avisos conhecidos,
  sem impacto funcional).
\item
  Testes funcionais em navegador (Chromium): \textbf{58 testes PASS / 0
  falhas}.
\item
  Scripts de validação e verificação de contraste criados e a funcionar.
\item
  Código organizado, modular e documentado.
\item
  Relatório técnico detalhado mantido em repositório.
\end{itemize}

\subsection{5. O QUE ESTÁ EM
ANDAMENTO}\label{o-que-estuxe1-em-andamento}

\begin{itemize}
\tightlist
\item
  \textbf{Migração de Páginas ao Novo Design (M4):} Aplicação do novo
  cabeçalho/rodapé e estrutura do Design System às restantes páginas do
  site.
\item
  \textbf{Estado actual:} Apenas a \textbf{Página Inicial (index.html)}
  foi migrada para o novo padrão. As restantes \textbf{14 páginas}
  encontram-se ainda com estrutura anterior, aguardando aplicação do
  novo layout.
\end{itemize}

\subsection{6. O QUE FALTA FAZER (Por
Concluir)}\label{o-que-falta-fazer-por-concluir}

{\def\LTcaptype{none} % do not increment counter
\begin{longtable}[]{@{}
  >{\raggedright\arraybackslash}p{(\linewidth - 4\tabcolsep) * \real{0.3333}}
  >{\raggedright\arraybackslash}p{(\linewidth - 4\tabcolsep) * \real{0.3333}}
  >{\raggedright\arraybackslash}p{(\linewidth - 4\tabcolsep) * \real{0.3333}}@{}}
\toprule\noalign{}
\begin{minipage}[b]{\linewidth}\raggedright
Item
\end{minipage} & \begin{minipage}[b]{\linewidth}\raggedright
Prioridade
\end{minipage} & \begin{minipage}[b]{\linewidth}\raggedright
Descrição
\end{minipage} \\
\midrule\noalign{}
\endhead
\bottomrule\noalign{}
\endlastfoot
\textbf{1. Migrar 14 páginas restantes} & \textbf{ALTA (Mais Urgente)} &
Aplicar header/footer + Design System em: Instituição, Órgãos Sociais,
Actividades, Notícias, Anúncios, Documentação, Despachantes, Parceiros,
Associar-se, Contactos, Galeria, Estatutos, Regulamentos. \\
\textbf{2. Garantir estrutura semântica} & ALTA & Aplicar
\texttt{\textless{}main\textgreater{}}, skip-link e estrutura acessível
em todas as páginas. \\
\textbf{3. Porting seletivo v3} (UX Leitura) & MÉDIA & Melhorias de
leitura/editorial (baixo risco, bom valor para o utilizador). \\
\textbf{4. Verificação Profissionais (v4)} & MÉDIA & Página de
verificação --- integrar com backend existente de forma controlada. \\
\textbf{5. Completar mapeamento PDFs} & BAIXA--MÉDIA & Reduzir avisos de
PDFs não declarados (sem impacto funcional). \\
\textbf{6. Acervo Fotográfico} & BAIXA & Integração futura do acervo
(865 imagens) na Galeria (fora do repositório por dimensão). \\
\end{longtable}
}

\subsection{7. CRONOGRAMA POR FASES (Dias
Úteis)}\label{cronograma-por-fases-dias-uxfateis}

{\def\LTcaptype{none} % do not increment counter
\begin{longtable}[]{@{}
  >{\raggedright\arraybackslash}p{(\linewidth - 4\tabcolsep) * \real{0.3333}}
  >{\raggedright\arraybackslash}p{(\linewidth - 4\tabcolsep) * \real{0.3333}}
  >{\raggedright\arraybackslash}p{(\linewidth - 4\tabcolsep) * \real{0.3333}}@{}}
\toprule\noalign{}
\begin{minipage}[b]{\linewidth}\raggedright
Fase
\end{minipage} & \begin{minipage}[b]{\linewidth}\raggedright
Actividades
\end{minipage} & \begin{minipage}[b]{\linewidth}\raggedright
Estado
\end{minipage} \\
\midrule\noalign{}
\endhead
\bottomrule\noalign{}
\endlastfoot
\textbf{Fase 1 --- Fundação} (22--23/09) & Identidade, logos, notícias,
parceiros & \textbf{Concluída} \\
\textbf{Fase 2 --- Levantamento} (01--03/10) & Acervo, dados, auditoria
técnica & \textbf{Concluída} \\
\textbf{Fase 3 --- Design System + Base Técnica} (06/10) & Tokens,
layout, componentes, datasets validados, testes & \textbf{Concluída} \\
\textbf{Fase 4 --- Migração de Páginas} (Pendente) & Aplicar novo design
às 14 páginas restantes & \textbf{Em Andamento (1/15)} \\
\textbf{Fase 5 --- Polimento e Homologação} (Após Fase 4) & Revisões,
testes finais, build/produção & \textbf{A Iniciar} \\
\end{longtable}
}

\emph{Obs.: Trabalho de Sábado foi consolidado em Sexta-feira. Sábado e
Domingo não considerados dias de trabalho.}

\subsection{8. INDICADOR DE CONCLUSÃO}\label{indicador-de-conclusuxe3o}

{\def\LTcaptype{none} % do not increment counter
\begin{longtable}[]{@{}
  >{\raggedright\arraybackslash}p{(\linewidth - 2\tabcolsep) * \real{0.5000}}
  >{\raggedright\arraybackslash}p{(\linewidth - 2\tabcolsep) * \real{0.5000}}@{}}
\toprule\noalign{}
\begin{minipage}[b]{\linewidth}\raggedright
Critério
\end{minipage} & \begin{minipage}[b]{\linewidth}\raggedright
Progresso
\end{minipage} \\
\midrule\noalign{}
\endhead
\bottomrule\noalign{}
\endlastfoot
\textbf{Estrutura Base (Design System + Dados + Componentes)} &
\textbf{100\% Concluído} \\
\textbf{Páginas Migradas para Novo Design} & \textbf{1 de 15 (6,7\%)} \\
\textbf{Progresso Global Estimado} & \textbf{\textasciitilde75\%
Concluído} \\
\end{longtable}
}

\subsection{9. RECOMENDAÇÃO}\label{recomendauxe7uxe3o}

\textbf{Prioridade Imediata:} Concluir \textbf{Fase 4 (Migração das 14
páginas restantes)}. Esta é a tarefa mais urgente para uniformizar todo
o site com o novo padrão visual e identidade da CDA, permitindo avançar
para polimento final e disponibilização.

\subsection{10. CONCLUSÃO}\label{conclusuxe3o}

O projecto encontra-se numa fase sólida e bem estruturada. A base
técnica, os dados e o Design System estão completamente definidos e
validados. O trabalho restante é essencialmente de \textbf{aplicação do
design já aprovado} às páginas existentes, o que permitirá concluir o
site de forma consistente e profissional.

\textbf{Responsável Técnico:} Equipa de Desenvolvimento\\
\textbf{Data de Emissão:} 07/10/2026
